\begin{table}[htbp]
\centering
\small
\caption{Sup-Wald structural break test for the volatility series}
\label{tab:break-test}
\begin{tabular}{lr}
\toprule
\textbf{Quantity} & \textbf{Value} \\
\midrule
Series tested & Yang-Zhang volatility (out-of-sample) \\
Candidate window & 1 Jan 2022 -- 1 Jan 2025 \\
Estimated break date & 8 February 2024 \\
Sup-Wald $F$-statistic & 1{,}809.05 \\
Critical value (10\%) & 5.74 \\
Critical value (5\%) & 7.17 \\
Critical value (1\%) & 10.17 \\
Pre-break observations & 10{,}098 \\
Post-break observations & 470 \\
Mean volatility, pre-break & 0.0134 \\
Mean volatility, post-break & 0.0313 \\
Post/pre ratio & $2.34\times$ \\
\bottomrule
\end{tabular}
\par\smallskip
{\footnotesize\textit{Note:} Single-break sup-Wald test
\parencite{andrews1993tests} applied to the out-of-sample Yang-Zhang
volatility series. The test fits a mean-shift regression at each candidate
date and reports the maximum $F$-statistic; the null hypothesis of no break
is rejected at the 1\% level. Critical values from \textcite{andrews1993tests},
Table~1 ($k = 1$, mean shift).}
\end{table}
